\documentclass[11pt,a4paper]{article}
\usepackage[margin=27mm]{geometry}
\usepackage[T1]{fontenc}
\usepackage{lmodern,microtype,amsmath,amssymb,amsthm,booktabs}
\usepackage[hidelinks]{hyperref}
\newcommand{\vol}{\operatorname{vol}}
\newcommand{\E}{\mathbb E}
\newcommand{\PP}{\mathbb P}
\newcommand{\CC}{\mathcal C}
\newtheorem{proposition}{Proposition}
\newtheorem{theorem}[proposition]{Theorem}
\newtheorem{lemma}[proposition]{Lemma}
\title{The dimension-nine threshold: an exact reduction\texorpdfstring{\\}{ }and minimality for constructions from paired seeds}
\author{Robin Gounder}
\date{10 October 2026\quad Research supplement}
\begin{document}
\maketitle
\begin{abstract}
The accompanying paper proves separation from the point-generated
product--join class in every dimension at least nine. We show that global
minimality of nine is equivalent to the unrestricted simplex bound in
dimension eight, which remains unproved here. We do prove that nine is
the first separation dimension within the paired two-upper-facet seed
family and for all constructions from those seeds by products and joins.
The proof uses a rational supporting gradient and log-concavity; a
separate positive-polynomial certificate checks the seed bound. Neither
argument reduces arbitrary convex bodies to these constructions.
\end{abstract}

\section{What global minimality requires}
For a full-dimensional convex body $K\subset\mathbb R^n$, put
\[
 R_n(K)=\frac{\vol_n(\Pi K)}{\vol_n(K)^{n-1}},\qquad
 g(n)=\frac{n^n}{n!},\qquad c_n=(n+1)g(n).
\]
Here $\Pi K$ is the projection body. Let $\CC$ be the class generated
from a point by Cartesian products, joins, and invertible affine maps
on the relevant affine hulls.
The exact finite classification in~\cite{finite} gives
$\sup_{K\in\CC_n}R_n(K)=c_n$ for $1\le n\le9$.
The accompanying paper~\cite{paper} supplies a strict witness in dimension
nine and separation in every larger dimension. The unrestricted simplex
bound in dimensions four through eight is left open in~\cite[Introduction]{feng}.

\begin{proposition}\label{prop:global}
Nine is the smallest dimension admitting separation from $\CC$ if and
only if
\begin{equation}\label{eq:remaining}
 R_8(K)\le c_8=\frac{131072}{35}
 \quad\text{for every full-dimensional convex body }K\subset\mathbb R^8.
\end{equation}
\end{proposition}
\begin{proof}
The inherited pyramid calculus~\cite[Sections 2--5]{calculus} gives,
for a $d$-dimensional convex body $K$ and $t\ge0$,
\[
 \frac{R_{d+t}(\mathcal P^tK)}{c_{d+t}}
 =\frac{R_d(K)}{g(d)}\frac{1+t a(K)}{d+t+1},
 \qquad a(K)\ge\frac1{d+1}.
\]
Here $a(K)$ is the lifted facet invariant of~\cite{calculus}, extended
by continuity to convex bodies.
The difference between consecutive terms on the right is
\[
 \frac{R_d(K)}{g(d)}
 \frac{(d+1)a(K)-1}{(d+t+1)(d+t+2)}\ge0.
\]
Thus any strict violation of the simplex bound in a dimension $d\le8$
lifts to a strict violation in dimension eight. If~\eqref{eq:remaining}
holds, no such lower-dimensional separation exists, and the known
dimension-nine witness proves minimality. Conversely, a violation
of~\eqref{eq:remaining} is already separation in dimension eight.
\end{proof}

By continuity of volume and the projection-body operator under Hausdorff
convergence, it would suffice to prove~\eqref{eq:remaining} for every
full-dimensional rational polytope in $\mathbb R^8$. This is still an
unbounded class; finite searches do not establish that assertion.

\section{A proved minimum for the paired family}
For $0\le u_1\le\cdots\le u_m<1$ and $n\ge2m$, define
\[
 K_n(u)=\{x\in\mathbb R^n:x_i\ge0,\
             \boldsymbol1\cdot x+|b\cdot x|\le1\},\qquad
 b=(-u_m,\ldots,-u_1,u_1,\ldots,u_m,0^{n-2m}).
\]
For $m=0$ this means the standard simplex. The family is inherited
from the two-facet constructions of~\cite{feng}; we make no priority
claim for the family or the dimension-nine phenomenon.

\begin{theorem}\label{thm:paired}
Every paired seed $K_n(u)$ in dimension $n\le8$ satisfies
$R_n(K_n(u))\le c_n$. The paired seed
$K_9(1/100,2/100,3/100,5/100)$ satisfies
\[
 \frac{R_9}{c_9}=\frac{20811283583}{20797677600}>1.
\]
Consequently nine is the first separation dimension within this paired
seed family, with zero coordinates and affine images allowed.
\end{theorem}

We first record an elementary consequence of log-concavity.
\begin{lemma}\label{lem:laplace}
If $\PP(Y>0)=1$, $Y$ has a log-concave survival function, and its mean
$M$ is finite, then $\E e^{-Y}\le(1+M)^{-1}$.
\end{lemma}
\begin{proof}
Write $S(t)=\PP(Y>t)$, so $S(0)=1$.
Convexity of $-\log S$ implies $S(t+s)\le S(t)S(s)$ for $s,t\ge0$.
Consequently
\[
 J(t):=\int_t^\infty S(s)\,ds\le M S(t).
\]
For $L=\int_0^\infty e^{-t}S(t)\,dt$, Tonelli's theorem gives
$\int_0^\infty e^{-t}J(t)\,dt=M-L\le ML$.
Hence $L\ge M/(1+M)$ and $\E e^{-Y}=1-L\le(1+M)^{-1}$.
\end{proof}

\begin{proof}[Proof of Theorem~\ref{thm:paired}]
It is enough to consider dimension eight: appending a zero tilt is a
pyramid, so Proposition~\ref{prop:global}'s monotonicity transfers the
eight-dimensional bound to smaller dimensions. Padding with zero pairs
then gives four parameters $0\le u_1\le\cdots\le u_4<1$.
Let $X_i$ be independent standard Laplace variables, each with density
$e^{-|x|}/2$, and write
\[
 M(u)=\E\left|\sum_{i=1}^4u_iX_i\right|,\qquad
 I(u)=8!\vol_8(K_8(u)).
\]
The homogeneous gauge integral gives
\[
 I(u)=\int_{\mathbb R_+^8}e^{-\boldsymbol1\cdot x-|b\cdot x|}\,dx
     =\E e^{-Y},\qquad Y=\left|\sum_{i=1}^4u_iX_i\right|,
\]
because the difference of two independent standard exponential
variables is standard Laplace. If $u\ne0$, the sum has a symmetric
log-concave density by Pr\'ekopa's convolution theorem~\cite[Theorem 7]{prekopa}.
The density of $Y$ on $[0,\infty)$ and its survival function are therefore
log-concave; the latter follows by integrating $f(t+s)$ over $s\ge0$
using~\cite[Theorem 6]{prekopa}. Lemma~\ref{lem:laplace} yields
$I(u)\le(1+M(u))^{-1}$.

The convex homogeneous function $M$ on $\mathbb R^4$ has, at $v=(1,2,4,8)$,
supporting gradient
\begin{equation}\label{eq:gradient}
 \nabla M(v)=\frac1{18225}(2749,5304,9624,15664).
\end{equation}
For a direct exact calculation, partial fractions of the characteristic
function $\prod_i(1+u_i^2t^2)^{-1}$ give
\[
 M(u)=\sum_{i=1}^4\frac{u_i^7}{\prod_{j\ne i}(u_i^2-u_j^2)}
 \quad(0<u_1<\cdots<u_4).
\]
Differentiating at $v$ gives~\eqref{eq:gradient}. Convexity and homogeneity
imply $M(u)\ge\nabla M(v)\cdot u$.
Writing $x_1=u_1$ and $x_j=u_j-u_{j-1}$ for $j\ge2$, we obtain
\begin{equation}\label{eq:gap}
 M(u)-\frac{u_1+3u_2+5u_3+7u_4}{9}
 \ge\frac{941x_1+217x_2+988x_3+1489x_4}{18225}\ge0.
\end{equation}
The seed formula in~\cite[Proposition 2]{paper} now gives, with
$w\cdot u=u_1+3u_2+5u_3+7u_4$,
\[
 \frac{R_8(K_8(u))}{c_8}
 =I(u)\left(1+\frac{w\cdot u}{9}\right)
 \le\frac{1+w\cdot u/9}{1+M(u)}\le1.
\]
For $u=0$ the body is a simplex; repeated and zero parameters also follow
by continuity, or directly from the integral argument. The displayed
rational value in dimension nine is checked by the exact seed formula
and the independent geometric certificate in~\cite{paper}. Since the
class maximum there is $c_9$, this is separation from $\CC_9$.
\end{proof}

\subsection{Closure under arbitrary products and joins}
Let $\mathcal F$ be the class generated by a point and all paired seeds
$K_n(u)$, using Cartesian products, joins, and invertible affine maps.
This allows arbitrary finite operation trees, rather than only pyramids
or the balanced recursion of the accompanying paper.

\begin{proposition}\label{prop:closure}
For $1\le n\le8$ and every $K\in\mathcal F_n$, one has $R_n(K)\le c_n$.
Thus nine is the first separation dimension within $\mathcal F$.
\end{proposition}
\begin{proof}
Use the inherited coordinates $H=1/a$ and $Q=aR/g(n)$ from~\cite{calculus}.
The universal invariant bounds give $1\le H\le n+1$. We prove by induction
on operation trees of output dimension at most eight that
\begin{equation}\label{eq:closed-envelope}
 Q(H+8-n)\le9.
\end{equation}
The formal point state is $n=0$, $H=Q=1$, which satisfies the bound.
For a seed of dimension $n\le8$, the left side divided by nine is
$R_8(\mathcal P^{8-n}K)/c_8$. These pyramids remain paired seeds, so
Theorem~\ref{thm:paired} proves the base cases.

For a join of factors with dimensions $r,s$, states $(h,q),(j,v)$, and
$n=r+s+1\le8$, the new state is $(h+j,qv)$. Put
$A=h-r-1\le0$ and $B=j-s-1\le0$. Then
\[
 (h+8-r)(j+8-s)=(9+A)(9+B)
 \ge9(9+A+B)=9(h+j+8-n).
\]
The two input bounds therefore imply the output bound.

For a product, assume $1\le r\le s$ and $n=r+s\le8$; products with a
point do not change a state. With $G=g(r)g(s)/g(n)$, the operation rules
give
\[
 H=\frac{nhj}{rj+sh},\qquad Q=qvG\frac{rj+sh}{n}.
\]
It suffices to verify
\begin{equation}\label{eq:product-envelope}
 (h+8-r)(j+8-s)
 -9G\left(hj+\frac{8-n}{n}(sh+rj)\right)\ge0.
\end{equation}
This is bilinear in $(h,j)$, so its minimum on
$[1,r+1]\times[1,s+1]$ occurs at a corner. For the sixteen pairs
$1\le r\le s$, $r+s\le8$, all 64 corner values are at least nine;
the checker recomputes each exact value. This proves product closure.

Finally, for $n\le8$, inequality~\eqref{eq:closed-envelope} implies
\[
 \frac{R_n(K)}{c_n}=\frac{HQ}{n+1}
 \le\frac{9H}{(n+1)(H+8-n)}\le1,
\]
where the last inequality is $(8-n)H\le(8-n)(n+1)$.
The known paired witness in dimension nine belongs to $\mathcal F$ and
exceeds the class maximum $c_9$.
\end{proof}

\section{Independent arithmetic check and remaining scope}
The file \texttt{check\_minimality.py} recomputes the supporting gradient,
the four margins in~\eqref{eq:gap}, the dimension-nine rational ratio,
and every corner obligation in~\eqref{eq:product-envelope}.
An unpruned enumeration of all operation states gives the simplex
maximum in dimensions one through nine, independently cross-checking
the previously published Pareto classification~\cite{finite}.

There is also a separate algebraic proof of the eight-dimensional family
bound. For distinct positive parameters, write
\[
 I(u)=\frac{N(u)}{D(u)},\qquad
 D(u)=\prod_i(1+u_i)\prod_{i<j}(u_i+u_j).
\]
The supplied integer numerator $N$ has 125 monomials. The checker
reconstructs $D$ and verifies the divided-difference formula for $I$
after multiplication by the Vandermonde polynomial. It then substitutes
$u_i=x_1+\cdots+x_i$ into
\[
 9D(u)-(9+w\cdot u)N(u).
\]
The resulting polynomial has 585 nonzero monomials, all with positive
integer coefficients, of total degrees seven through ten. This proves
the bound throughout the ordered positive parameter region; continuity
handles its boundary. The checker uses standard-library integer and
rational arithmetic, with no computer-algebra system or optimizer.

Run the following commands from this supplement's directory:
\begin{quote}\small\ttfamily
python -B check\_minimality.py -{}-check-report\\
python -B -m unittest -v test\_minimality.py\\
python -B -O -m unittest -v test\_minimality.py
\end{quote}
The tests reject altered numerators and denominators, false or missing
gap coefficients, duplicate monomials, and invalid exponents. The
analytic arguments are supplied above, rather than inferred from finite
checks, and have not been proof-assistant formalized or externally
peer reviewed.

Theorem~\ref{thm:paired} and Proposition~\ref{prop:closure} exclude the paired
family and its complete product--join closure as sources of a
lower-dimensional witness. General asymmetric seeds and arbitrary convex
bodies need not belong to this generated class. The results do not
prove~\eqref{eq:remaining}, nor imply that every maximizer belongs to
$\mathcal F$. Global minimality is therefore open in this work. The next
substantive requirement is a bound covering all eight-dimensional convex
bodies, or an exact counterexample refuting it.

\section*{Generative AI disclosure}
OpenAI Codex agents contributed substantially to source discovery,
literature synthesis, and mathematical exposition.

\section*{Data and code availability}
This note, its source, polynomial certificate, checker, and exact-arithmetic
report are supplied in the accompanying reproduction repository~\cite{paper},
under \texttt{research/minimality/}. The original paper and its versioned
release remain separate from this research supplement. The supplied
checks establish the stated arithmetic obligations; they do not claim
an unrestricted nonexistence proof or global bibliographic priority.

\begin{thebibliography}{9}
\bibitem{feng}
Y.~Feng, S.~Hu, W.~Liu, and L.~Xu,
\emph{On the reverse projection inequality}, author-deposited preprint,
MathSciDoc:2608.23002, PDF dated 26 August 2026.
\href{https://archive.ymsc.tsinghua.edu.cn/pacm_paperurl/20260821091122099811781}{Primary archive record}.

\bibitem{paper}
R.~Gounder, \emph{Beyond point-generated product--join constructions
for projection-body volume}, working research draft, version 1.1.0, 2026.
\href{https://github.com/Kohze/projection-body-separation/releases/tag/v1.1.0}{Versioned manuscript and certificates};
\href{https://github.com/Kohze/projection-body-separation/tree/main/research/minimality}{Minimality research supplement}.

\bibitem{finite}
mxym (repository account), \emph{An exact Pareto dynamic program for
product--join projection geometry: sharp extrema through dimension 48},
research note, 8 October 2026.
\href{https://github.com/mxym/math/blob/55ba5a7be927e3e7d2dbfd8d054db91b63d8d755/notes/exact-product-join-finite-optima/paper.md}{Pinned theorem and proof}.

\bibitem{calculus}
mxym (repository account), \emph{Projection-volume calculus for joins
and Cartesian products}, research manuscript, version 2, 7 October 2026.
\href{https://github.com/mxym/math/blob/2fc3f25179cf1128c4bccbde65fd3a07c093ad6d/preprints/005-simplex-product-optimum/v2/paper.md}{Pinned manuscript}.

\bibitem{prekopa}
A.~Pr\'ekopa, \emph{On logarithmic concave measures and functions},
Acta Scientiarum Mathematicarum (Szeged) \textbf{34} (1973), 335--343.
\href{https://acta.bibl.u-szeged.hu/14411/1/math_034_335-343.pdf}{Published full text}.
\end{thebibliography}
\end{document}
